% p401_1; source PDF page 401; manually verified editable table.
\begin{tikzpicture}[x=1cm,y=-1cm,inner sep=0pt,outer sep=0pt]
\path[use as bounding box] (0,0) rectangle (16.000,4.495);
\draw[line width=.4pt] (0.000,0.000) rectangle (4.200,1.215);
\node[align=center,text width=3.940cm,font={\linespread{1}\fontsize{10}{13.50}\selectfont\bfseries}] at (2.100,0.608) {非国家工作人员（甲）};
\draw[line width=.4pt] (4.200,0.000) rectangle (7.750,1.215);
\node[align=center,text width=3.290cm,font={\linespread{1}\fontsize{10}{13.50}\selectfont\bfseries}] at (5.975,0.608) {国家工作人员（乙）};
\draw[line width=.4pt] (7.750,0.000) rectangle (16.000,1.215);
\node[align=center,text width=7.990cm,font={\linespread{1}\fontsize{10}{13.50}\selectfont\bfseries}] at (11.875,0.608) {结论};
\draw[line width=.4pt] (0.000,1.215) rectangle (4.200,2.730);
\node[align=center,text width=3.940cm,font={\linespread{1}\fontsize{10}{13.50}\selectfont}] at (2.100,1.973) {实行行为};
\draw[line width=.4pt] (4.200,1.215) rectangle (7.750,2.730);
\node[align=left,text width=3.290cm,font={\linespread{1}\fontsize{10}{13.50}\selectfont}] at (5.975,1.973) {实行行为};
\draw[line width=.4pt] (7.750,1.215) rectangle (16.000,2.730);
\node[align=left,text width=7.990cm,font={\linespread{1}\fontsize{10}{13.50}\selectfont}] at (11.875,1.973) {甲乙既构成财产犯罪的共同犯罪，也构成贪污罪的共同犯罪，想象竞合，择一重罪论处（例2）};
\draw[line width=.4pt] (0.000,2.730) rectangle (4.200,4.495);
\node[align=center,text width=3.940cm,font={\linespread{1}\fontsize{10}{13.50}\selectfont}] at (2.100,3.612) {实行行为};
\draw[line width=.4pt] (4.200,2.730) rectangle (7.750,4.495);
\node[align=left,text width=3.290cm,font={\linespread{1}\fontsize{10}{13.50}\selectfont}] at (5.975,3.612) {帮助行为};
\draw[line width=.4pt] (7.750,2.730) rectangle (16.000,4.495);
\node[align=left,text width=7.990cm,font={\linespread{1}\fontsize{10}{13.50}\selectfont}] at (11.875,3.612) {甲构成财产犯罪（实行犯），乙构成财产犯罪（帮助犯）。乙不构成贪污罪（例3）};
\end{tikzpicture}
